\chapter{Frontend Error Boundaries}

Not every frontend error is an API error. A malformed prop, a
\texttt{undefined.map()}, or a third-party component throwing during render
will crash the entire React tree with a blank white screen unless something
catches it. That "something" is an Error Boundary.

\section{Render Errors vs. API Errors}

\begin{warnbox}[WARNING]
Error Boundaries only catch errors thrown during \textbf{rendering},
\textbf{lifecycle methods}, and \textbf{constructors} of the component tree
below them. They do \emph{not} catch errors inside event handlers, async
code (\texttt{setTimeout}, promises), or \texttt{fetch}/\texttt{axios}
calls. A failed API request is handled by the try/catch and \texttt{error}
state pattern from Chapter~33 --- Error Boundaries are a separate,
complementary safety net for unexpected render-time crashes.
\end{warnbox}

\begin{diagrambox}[What Each Mechanism Catches]
\begin{verbatim}
 Error type                  | Caught by
 ----------------------------|----------------------------
 Component throws while      | Error Boundary
 rendering (bad prop, null   |
 access, etc.)                |
 ----------------------------|----------------------------
 API call fails (network,    | try/catch + error state
 4xx/5xx response)            | (Chapter 33)
 ----------------------------|----------------------------
 Error thrown inside an      | Local try/catch in the
 onClick/onSubmit handler     | handler itself
\end{verbatim}
\end{diagrambox}

\section{Building an ErrorBoundary}

Error Boundaries must currently be class components --- there is no hook
equivalent for \texttt{getDerivedStateFromError} or
\texttt{componentDidCatch}.

\begin{lstlisting}[language=JavaScript]
// components/ErrorBoundary.jsx
import { Component } from "react";

class ErrorBoundary extends Component {
  constructor(props) {
    super(props);
    this.state = { hasError: false, error: null };
  }

  static getDerivedStateFromError(error) {
    // Runs during the render phase; return new state to trigger the fallback UI.
    return { hasError: true, error };
  }

  componentDidCatch(error, errorInfo) {
    // Runs after render; safe place to log to a monitoring service.
    console.error("ErrorBoundary caught:", error, errorInfo);
  }

  handleReset = () => {
    this.setState({ hasError: false, error: null });
  };

  render() {
    if (this.state.hasError) {
      return (
        <div className="p-6 text-center">
          <h2 className="text-lg font-semibold text-red-600 mb-2">
            Something went wrong.
          </h2>
          <p className="text-sm text-gray-600 mb-4">
            {this.state.error?.message || "An unexpected error occurred."}
          </p>
          <button
            onClick={this.handleReset}
            className="px-4 py-2 bg-blue-600 text-white rounded"
          >
            Try Again
          </button>
        </div>
      );
    }

    return this.props.children;
  }
}

export default ErrorBoundary;
\end{lstlisting}

\begin{notebox}[NOTE]
\texttt{getDerivedStateFromError} must be a static method returning the new
state; it cannot perform side effects like logging. \texttt{componentDidCatch}
is where you send the error to a logging service (Sentry, LogRocket, or your
own endpoint) since it runs after the render phase completes.
\end{notebox}

\section{Wrapping the App}

Wrap boundaries around independent sections rather than one boundary around
the entire app --- that way a crash in one section (say, a chart widget)
doesn't take down the navbar and the rest of the page with it.

\begin{lstlisting}[language=JavaScript]
// App.jsx
import { BrowserRouter, Routes, Route } from "react-router-dom";
import ErrorBoundary from "./components/ErrorBoundary";
import Navbar from "./components/Navbar";
import TaskListPage from "./pages/TaskListPage";
import ProfilePage from "./pages/ProfilePage";

function App() {
  return (
    <BrowserRouter>
      {/* Navbar has its own boundary so a nav crash doesn't hide the rest of the page */}
      <ErrorBoundary>
        <Navbar />
      </ErrorBoundary>

      <ErrorBoundary>
        <Routes>
          <Route path="/tasks" element={<TaskListPage />} />
          <Route path="/profile" element={<ProfilePage />} />
        </Routes>
      </ErrorBoundary>
    </BrowserRouter>
  );
}

export default App;
\end{lstlisting}

\begin{tipbox}[TIP]
Granular boundaries (per route, or per risky widget) contain the damage of a
crash to just that section. A single app-wide boundary is easier to set up
but turns any unrelated bug into a full-page failure with no navigation left
to escape from.
\end{tipbox}

\section{Error Boundaries Do Not Replace API Error Handling}

Wrapping a component in an \texttt{ErrorBoundary} does nothing for a failed
\texttt{axios} call inside its \texttt{useEffect} --- that promise
rejection happens outside React's render cycle, so the boundary never sees
it. Every data-fetching component still needs its own \texttt{loading} /
\texttt{error} / retry handling as covered in Chapter~33. Error Boundaries
exist purely to stop unexpected render crashes from blanking the whole
screen; they are a last line of defense, not a substitute for handling
expected failure modes like a down API or a 500 response.

\whatwhyhow
{A class-based ErrorBoundary uses \texttt{getDerivedStateFromError} and \texttt{componentDidCatch} to catch render/lifecycle errors in its child tree and show a fallback UI instead of a blank screen.}
{Uncaught render errors crash the entire React tree with no recovery path; API/async errors need separate handling because boundaries cannot see them.}
{Wrap independent sections of the app (navbar, route content, risky widgets) in their own \texttt{<ErrorBoundary>}, and keep using try/catch plus loading/error state (Chapter~33) for every API call.}
